\exercisegroup

\begin{enumerate}
\def\labelenumi{\arabic{enumi})}
\tightlist
\item
  设 E 与 F 是两个满足第一可数公理的无穷维希尔伯特空间，\((a_{n})\) 是 E 的一个标准正交基，\((b_{n})\) 是 F 的一个标准正交基。
\end{enumerate}

\begin{enumerate}
\def\labelenumi{\alph{enumi})}
\tightlist
\item
  设 u 是从 E 到 F 中的连续线性映射；设 \(u(a_{n}) = \sum_{m} \alpha_{mn} b_{m}\)。证明
\[
\sum_ {n} | \alpha_ {m n} | ^ {2} \leqslant \| u \| ^ {2}
\]

\[
\sum_ {m} | \alpha_ {m n} | ^ {2} \leqslant \| u \| ^ {2}
\]
对每个 \(m\) 与 \(n\) 成立。

\begin{enumerate}
\def\labelenumi{\alph{enumi})}
\setcounter{enumi}{1}
\tightlist
\item
  给出一个二重序列 \((\alpha_{mn})\) 的例子，使得 \(\sum_{m} |\alpha_{mn}|^{2} \leqslant 1\) 对所有 \(n\) 成立，且 \(\sum_{n} |\alpha_{mn}|^{2} \leqslant 1\) 对所有 \(m\) 成立，但不存在任何从 E 到 F 中的连续线性映射 \(u\)，使得 \(\langle u(a_n)|b_m\rangle = \alpha_{mn}\) 对每对整数 \((m, n)\) 成立。（证明：若 I \(\subset\) N 是由 \(p\) 个整数组成的集，且 V \(_p\)（相应地 W \(_p\)）是 E（相应地 F）中由 \(a_n\)（相应地 \(b_n\)，均满足 \(n \in I\)）生成的子空间，则存在一个线性映射 \(u_p\)（从 V \(_p\) 到 W \(_p\) 上的），使得 \(\langle u_p(a_n)|b_m\rangle = \frac{1}{\sqrt{p}}\) 对 \(m \in I\) 与 \(n \in I\) 成立，且 \(\|u_p\| \geqslant \sqrt{p}\)。）
\end{enumerate}

\begin{enumerate}
\def\labelenumi{\arabic{enumi})}
\setcounter{enumi}{1}
\tightlist
\item
  设 \(A = (\alpha_{mn})_{(m,n)\in \mathbf{N}\times \mathbf{N}}\) 是一个复数的二重序列，我们也称它为无穷矩阵。对每个点 \(x = (x_{n})\)（直和空间 \(\mathbf{C}^{(\mathbf{N})}\) 中的），和式 \(y_{m} = \sum_{n}\alpha_{mn}x_{n}\) 都有定义；令 \(A.x\) 表示点 \((y_{m})\)，它属于乘积空间 \(\mathbf{C}^{\mathbf{N}}\)，于是 \(x\mapsto A.x\) 是从 \(\mathbf{C}^{(\mathbf{N})}\) 到 \(\mathbf{C}^{\mathbf{N}}\) 中的一个线性映射，且从 \(\mathbf{C}^{(\mathbf{N})}\) 到 \(\mathbf{C}^{\mathbf{N}}\) 中的每个线性映射都具有这一形式。用 \(E_{n}\) 表示 \(\mathbf{C}^{(\mathbf{N})}\) 中由典范基前 \(n\) 个向量生成的子空间，用 \(P_{n}\) 表示从 \(\mathbf{C}^{\mathbf{N}}\) 到 \(E_{n}\) 上的典范投影；当赋给 \(E_{n}\) 由空间 \(\ell_{\mathbf{C}}^{2}(\mathbf{N})\) 的范数诱导的范数时，\(\| u\|\) 表示线性映射 \(u\)（从有限维希尔伯特空间 \(E_{n}\) 到其自身的）的范数。
\end{enumerate}

\begin{enumerate}
\def\labelenumi{\alph{enumi})}
\item
  为使 \(\mathbf{C}^{(\mathbf{N})}\) 在映射 \(x\mapsto A.x\) 下的像含于 \(\ell_{\mathbf{C}}^2 (\mathbf{N})\)，且该映射可延拓为从 \(\ell_{\mathbf{C}}^2 (\mathbf{N})\) 到其自身的连续线性映射，必要且充分的是诸范数 \(\| P_nAP_n\|\) 有界。这蕴含 \(A\) 的行与列都属于 \(\ell_{\mathbf{C}}^2 (\mathbf{N})\)（习题 1）。
\item
  用 \(A^*\) 表示无穷矩阵 \((\alpha_{mn}')\)，其中 \(\alpha_{mn}' = \overline{\alpha}_{nm}\)。若 \(A\) 的列都属于 \(\ell_{\mathbf{C}}^2(\mathbf{N})\)（换言之，若 \(x \mapsto A.x\) 把 \(\mathbf{C}^{(\mathbf{N})}\) 映到 \(\ell_{\mathbf{C}}^2(\mathbf{N})\) 中），则级数 \(\beta_{mn} = \sum_{p} \overline{\alpha}_{pm} \alpha_{pn}\)
\end{enumerate}

绝对收敛，此时我们置 \(A^{*}A = (\beta_{mn})\)。证明：为使 \(x \mapsto A.x\) 可延拓为一个连续线性映射 \(u\)（从 \(\ell_{\mathbf{C}}^{2}(\mathbf{N})\) 到其自身的），必要且充分的是诸范数 \(\| P_n(A^*A)P_n\|\) 有界（我们有 \(\langle P_n(A^*A)P_n.x|x\rangle = \| AP_n .x\|^2\)，对所有 \(x\in E_n\) 成立）。于是 \(x\mapsto A^{*}A.x\) 延拓为一个正埃尔米特映射 \(u^{*}u\)（从 \(\ell_{\mathbf{C}}^{2}(\mathbf{N})\) 到其自身的）。

\begin{enumerate}
\def\labelenumi{\alph{enumi})}
\setcounter{enumi}{2}
\item
  对两个无穷矩阵 \(X = (\xi_{mn})\)、\(Y = (\eta_{mn})\)，我们称乘积 \(XY\) 有定义，如果级数 \(\zeta_{mn} = \sum_{p} \xi_{mp} \eta_{pn}\) 绝对收敛，此时我们置 \(XY = (\zeta_{mn})\)。我们称幂 \(X^k\)（\(k\) 为整数 \(> 1\)）有定义，如果 \(X^{k-1}\) 与 \(X^{k-1}X\) 都有定义，此时我们置 \(X^k = X^{k-1}X\)；在这一情形，\(X^pX^q = X^k\) 对每对整数 \(p, q\)（满足 \(p + q = k\)）成立。设 \(A\) 是一个列都属于 \(\ell_{\mathbf{C}}^2(\mathbf{N})\) 的无穷矩阵，且乘积 \((A^*A)^2\) 有定义，证明：对所有 \(x \in E_n\)，我们有 \(\langle (P_nA^*AP_n)^2.x|x\rangle \leqslant \langle P_n(A^*A)^2P_n.x|x\rangle\)，并由此推出 \(\| P_nA^*AP_n\|^2 \leqslant \| P_n(A^*A)^2P_n\|\)。
\item
  为使无穷矩阵 A 满足：\(\mathbf{C}^{(\mathbf{N})}\) 在 \(x\mapsto A.x\) 下的像含于 \(\ell_{\mathbf{C}}^2 (\mathbf{N})\)，且 \(x\mapsto A.x\) 可延拓为从 \(\ell_{\mathbf{C}}^2 (\mathbf{N})\) 到其自身的连续线性映射，必要且充分的是下列三个条件都被满足：
\end{enumerate}

\begin{enumerate}
\def\labelenumi{(\roman{enumi})}
\item
  \(A\) 的行与列都属于 \(\ell_{\mathbf{C}}^{2}(\mathbf{N})\)；
\item
  幂 \((A^{*}A)^{k}\) 对每个整数 \(k > 1\) 都有定义；
\item
  我们有
\end{enumerate}

\[
\sup _ {n} \left(\sup _ {m} | ((A ^ {*} A) _ {m m} ^ {n}) ^ {1 / n} |\right) <   + \infty
\]

其中 \((A^{*}A)_{mm}^{n}\) 表示以 \((m,m)\) 为指标的项，它属于矩阵 \((A^{*}A)^{n}\)。（注意，若 \(C\) 是 \(\mathrm{E}_n\) 的一个正自伴自同态关于 \(\mathrm{E}_n\) 的典范基的矩阵，则 \(\| C\| \leqslant n\sup_{1\leqslant i\leqslant n}|C_{ii}|\)（考察 \(C\) 的迹并将 \(C\) 对角化）。利用在 \(c\) ) 中证明的不等式，证明

\[
\left\| P _ {n} A ^ {*} A P _ {n} \right\| \leqslant n ^ {2 - k} \sup _ {1 \leqslant i \leqslant n} \left| ((A ^ {*} A) _ {i i} ^ {2 k}) \right| ^ {2 - k}
\]

对每个整数 \(k > 1\) 成立，只要条件 (i)、(ii) 与 (iii) 都被满足。）

¶ 3) 设 \((a_{ij})_{(i,j)\in\mathbb{I}\times\mathbb{I}}\) 是一个可数无穷的复数二重族。假设存在两个数 \(\beta > 0\)、\(\gamma > 0\)，以及一个由 > 0 的数组成的族 \((p_i)_{i\in\mathbb{I}}\)，满足关系式

\[
\sum_ {i} p _ {i} | a _ {i j} | \leqslant \beta p _ {j}, \quad \sum_ {j} | a _ {i j} | p _ {j} \leqslant \gamma p _ {i}\tag{*}
\]

对所有 \(i, j\)（属于 \(\mathbb{I}\)）成立。
a) 证明存在一个连续自同态 \(u\)（\(\ell_{\mathbf{C}}^2 (\mathrm{I})\) 上的），其范数 \(\leqslant (\beta \gamma)^{1 / 2}\)，使得对所有 \(x = (x_{i})_{i\in \mathbf{I}}\)（属于 \(\ell_{\mathbf{C}}^2 (\mathrm{I})\)），有 \(u(x) = y\)，其中 \(y = (y_{i})\) 由 \(y_{i} = \sum_{j}a_{ij}x_{j}\) 给出（对 \(x = (x_{i})\) 与 \(y = (y_{i})\)（\(\mathbf{C}^{(\mathbb{I})}\) 中的），置 \(v_{ij} = |x_i| (p_j|a_{ij}| / p_i)^{1 / 2}\)、\(w_{ij} = |y_j|(p_i|a_{ij}| / p_j)^{1 / 2}\)，并给出 \(\sum_{i,j}v_{ij}w_{ij}\) 的一个上界）。

\begin{enumerate}
\def\labelenumi{\alph{enumi})}
\setcounter{enumi}{1}
\tightlist
\item
  取 I 为所有整数 \(\geqslant 1\) 组成的集，并设 \(a_{ij} = (i + j)^{-1}\)。证明当 \(p_i = i^{-1/2}\) 且 \(\beta = \gamma = \pi\) 时条件 (*) 被满足（这些是可能取到的最佳常数）（把 (*) 中的级数与一个积分相比较）。以类似的方式处理 I = N 且 \(a_{ij} = (i + j + 1)^{-1}\) 的情形（「Hilbert 矩阵」）。
\end{enumerate}

* c) 设 \(\mathcal{H} = \mathrm{H}^2 (\mathbf{D})\) 是 Hardy 空间，它由所有满足下列条件的函数 \(f(z) = \sum_{n=0}^{\infty} a_n z^n\) 组成：在开圆盘 \(\mathbf{D}:|z| < 1\) 中解析，且 \(\| f\|^2 = \sum_n |a_n|^2 < +\infty\)；于是 \(\| f\|\) 是 \(\mathcal{H}\) 上的一个范数，对于它 \(\mathcal{H}\) 同构于 \(\ell_{\mathbf{C}}^2 (\mathbf{N})\)。给定两个函数 \(f, g\)（属于 \(\mathcal{H}\)），证明函数 \(t \mapsto f(t) g(t)\)（实变量 \(t\) 的）在 \([0,1]\) 上关于 Lebesgue 测度可积，且公式 \(\mathbf{B}(f,g) = \int_0^1 f(t) g(t) dt\) 在 \(\mathcal{H} \times \mathcal{H}\) 上定义一个连续双线性型（考察 \(\mathbf{B}(f,f) = \int_0^1 f(t)^2 dt\)，其中函数 \(f \in \mathcal{H}\) 形如 \(f(z) = \sum_{n=0}^{\mathbf{N}} a_n z^n\)，且 \(a_n \geqslant 0\)（对 \(0 \leqslant n \leqslant \mathbf{N}\)）；用 Cauchy 定理建立关系式

\[
\int_ {- 1} ^ {1} f (t) ^ {2} d t = - i \int_ {0} ^ {\pi} f (e ^ {i \theta}) ^ {2} e ^ {i \theta} d \theta
\]

由此得到 \(\mathbf{B}(f, f) \leqslant \frac{1}{2} \int_{-\pi}^{\pi} |f(e^{i\theta})|^2 d\theta = \pi \| f\|^2\)。用这一方法，得到 \(b\) ) 的结果：Hilbert 矩阵定义了一个范数 \(\leqslant \pi\) 的自同态（在希尔伯特空间 \(\ell_{\mathbf{C}}^2(\mathbf{N})\) 上）。*

\begin{enumerate}
\def\labelenumi{\arabic{enumi})}
\setcounter{enumi}{3}
\tightlist
\item
  设 \(\mathbf{E}\) 是一个维数为 \(d\) 的复希尔伯特空间。
\end{enumerate}

\begin{enumerate}
\def\labelenumi{\alph{enumi})}
\item
  对每个 \(u \geqslant 0\)（\(\mathcal{L}(\mathrm{E})\) 中的，V, p.~45），证明存在唯一的 \(v \geqslant 0\)（\(\mathcal{L}(\mathrm{E})\) 中的）使 \(u = v^2\)（把 \(u\) 对角化）；我们记 \(v = u^{1/2}\)。
\item
  对每个 \(u \in \mathcal{L}(\mathrm{E})\)，置 \(\mathrm{abs}(u) = (u^{*}u)^{1/2}\)。证明 u 与 \(\mathrm{abs}(u)\) 有相同的范数，且我们有 \(\mathrm{abs}(\symbf{\Lambda}^{n}(u)) = \symbf{\Lambda}^{n}(\mathrm{abs}(u))\)（对每个整数 \(n \leqslant d\)）。
\item
  设 \(s_{1}(u) \geqslant s_{2}(u) \geqslant \ldots \geqslant s_{d}(u) \geqslant 0\) 是 \(\mathrm{abs}(u)\) 的按重数计算的特征值序列。证明 \(\|u\| = s_{1}(u)\)，且对每个整数 \(n \leqslant d\)，\(\| \symbf{\Lambda}^{n}(u)\| = s_{1}(u) s_{2}(u) \ldots s_{n}(u)\)。为使 \(\| \symbf{\Lambda}^{n}(u)\| = \| u\|^{n}\)（对所有满足 \(1 \leqslant n \leqslant d\) 的 n），必要且充分的是 \(u^{*}u\) 是一个位似，换言之，u 是一个酉算子的标量倍。
\end{enumerate}

\begin{enumerate}
\def\labelenumi{\arabic{enumi})}
\setcounter{enumi}{4}
\tightlist
\item
  设 E、F 是两个希尔伯特空间，u 是从 E 到 F 中的连续线性映射。用 \(\ell(u)\) 表示所有 \(x \in E\)（满足 \(\|u(x)\| = \|u\|. \|x\|\)）组成的集。
\end{enumerate}

\begin{enumerate}
\def\labelenumi{\alph{enumi})}
\item
  证明 \(\ell(u)\) 是 E 的闭向量子空间，它是 \(u^{*}u - \| u\|^{2}1_{\mathrm{E}}\) 的核，且与 \(u\) 的核正交。
\item
  证明 \(u\) 在 \(\ell(u)\) 上的限制是从 \(\ell(u)\) 到 \(\ell(u^*)\) 上的双射，其逆双射是 \(\|u\|^{-2}.u^*\) 在 \(\ell(u^*)\) 上的限制；此外，正交补 \((\ell(u))^{\circ}\) 在 \(u\) 下的像含于 \((\ell(u^*))^{\circ}\)。若 \(u_1\) 是 \(u\) 在 \((\ell(u))^{\circ}\) 上的限制（视为从 \((\ell(u))^{\circ}\) 到 \((\ell(u^*))^{\circ}\) 中的映射），则其伴随 \(u_1^*\) 是 \(u^*\) 在 \((\ell(u^*))^{\circ}\) 上的限制；若 \(\ell(u) \neq E\)，用 \(\langle u\rangle\)（称为 \(u\) 的次范数）表示范数 \(\|u_1\|\)；若 \(\ell(u) = E\)，则置 \(\langle u\rangle = 0\)。于是 \(\langle u^*\rangle = \langle u\rangle\)。
\end{enumerate}

\begin{enumerate}
\def\labelenumi{\arabic{enumi})}
\setcounter{enumi}{5}
\tightlist
\item
  设 E 是希尔伯特空间，M、N 是 E 的两个闭子空间，\(M^{\circ}\)、\(N^{\circ}\) 是它们各自的正交补；用 \(p_{M}\)、\(p_{N}\) 分别表示到 M 与 N 上的正交投影，使 \(1_{E} - p_{M}\)、\(1_{E} - p_{N}\) 分别是到 \(M^{\circ}\) 与 \(N^{\circ}\) 上的正交投影。置 \(u_{\mathrm{NM}} = (1_{\mathrm{E}} - p_{\mathrm{N}}) p_{\mathrm{M}}\) 与 \(\delta(\mathbf{M}, \mathbf{N}) = \|u_{\mathrm{NM}}\|\)；我们有 \(\delta(\mathbf{M}, \mathbf{N}) = \delta(\mathbf{N}^{\circ}, \mathbf{M}^{\circ}) \leqslant 1\)；关系 \(\delta(\mathbf{M}, \mathbf{N}) < 1\) 蕴含 \(M \cap N^{\circ} = \{0\}\)。
\end{enumerate}

\begin{enumerate}
\def\labelenumi{\alph{enumi})}
\item
  用 \(\tilde{\mathbf{M}}\) 表示 \(\mathbf{M}\) 中 \(\mathbf{M} \cap \mathbf{N}^{\circ}\) 的正交补，并设 \(\varepsilon(\mathbf{M}, \mathbf{N}) = \delta(\tilde{\mathbf{M}}, \mathbf{N})\)。证明（用习题 5 的记号）\(\ell(u_{\mathrm{NM}}) = \mathbf{M} \cap \mathbf{N}^{\circ}\)，并由此推出 \(\varepsilon(\mathbf{M}, \mathbf{N}) = \langle u_{\mathrm{NM}} \rangle \leqslant \delta(\mathbf{M}, \mathbf{N})\)；此外，若 \(\mathbf{M} \cap \mathbf{N}^{\circ} = \{0\}\)（特别地，若 \(\delta(\mathbf{M}, \mathbf{N}) < 1\)），则 \(\varepsilon(\mathbf{M}, \mathbf{N}) = \delta(\mathbf{M}, \mathbf{N})\)。
\item
  对从 E 到其自身的连续线性映射 \(u\)，我们称 \(u\) 的余范数为数 \(c(u) = \inf \|u(x)\|/\|x\|\)，其中 \(x\) 取遍所有 \(\neq 0\) 且与 \(u^{-1}(0)\) 正交的向量组成的集（若 \(u = 0\)，置 \(c(u) = 1\)）。为使 \(u(\mathrm{E})\) 在 E 中是闭的，必要且充分的是 \(c(u) > 0\)（I, p.~17, th. 1）。我们有 \(c(u^{*}) = c(u)\)。
\item
  设 \(v_{\mathbf{NM}} = p_{\mathbf{N}}p_{\mathbf{M}}\)。证明
\end{enumerate}

\[
\varepsilon (\mathbf {M}, \mathbf {N}) ^ {2} + c (v _ {\mathbf {N M}}) ^ {2} = 1
\]

（注意 \(\|u_{\mathrm{NM}}(x)\|^{2} + \|v_{\mathrm{NM}}(x)\|^{2} = \|p_{\mathrm{M}} \cdot x\|^{2}\)，且 \(v_{NM}\) 的核是 \(M^{\circ} + (M \cap N^{\circ})\)，并由此推出 \(\langle u_{NM}\rangle^{2} \leqslant 1 - c(v_{\mathrm{NM}})^{2}\)。）
\item
  由 b) 与 c) 推出 \(\varepsilon(\mathrm{N}, \mathrm{M}) = \varepsilon(\mathrm{M}, \mathrm{N})\)，并利用 a) 推出 \(\varepsilon(\mathrm{M}^{\circ}, \mathrm{N}^{\circ}) = \varepsilon(\mathrm{M}, \mathrm{N})\)。
\item
  置 \(g(\mathbf{M}, \mathbf{N}) = \|p_{\mathbf{M}} - p_{\mathbf{N}}\|\)（cf.~V, p.~64, 习题 17）。证明

\[
g (\mathbf {M}, \mathbf {N}) = \sup (\delta (\mathbf {M}, \mathbf {N}), \delta (\mathbf {N}, \mathbf {M}))
\]

（注意 \(p_{\mathbf{M}} - p_{\mathbf{N}} = (1_{\mathrm{E}} - p_{\mathbf{N}})p_{\mathbf{M}} - p_{\mathbf{N}}(1_{\mathrm{E}} - p_{\mathbf{M}})\)）；由此推出 \(\varepsilon (\mathbf{M},\mathbf{N})\leqslant g(\mathbf{M},\mathbf{N})\)。若 \(\mathbf{M}\cap \mathbf{N}^{\circ} = \mathbf{N}\cap \mathbf{M}^{\circ} = \{0\}\)，则我们有关系式 \(\varepsilon (\mathbf{M},\mathbf{N}) = \delta (\mathbf{M},\mathbf{N}) = \delta (\mathbf{N},\mathbf{M}) = g(\mathbf{M},\mathbf{N})\)。若 \(g(\mathbf{M},\mathbf{N}) < 1\)，则 \(\mathbf{M}\cap \mathbf{N}^{\circ} = \mathbf{N}\cap \mathbf{M}^{\circ} = \{0\}\)。

\end{enumerate}

\begin{enumerate}
\def\labelenumi{\alph{enumi})}
\setcounter{enumi}{5}
\tightlist
\item
  设 \(Q_{M}\)、\(Q_{N}\) 是 E 中两个连续投影，它们的像分别为 M 与 N；给出关系式 \(g(\mathrm{M}, \mathrm{N}) \leqslant \|Q_{\mathrm{M}} - Q_{\mathrm{N}}\|\) 的另一个证明（V, p.~64, exerc. 17）。（注意，对所有 \(x \in E\)，我们有 \(\|(1_{E} - Q_{M}).x\|^{2} + \|Q_{M}^{*}.x\|^{2} = \|x\|^{2} + \|(Q_{\mathrm{M}} - Q_{\mathrm{M}}^{*}).x\|^{2}\)，把这一关系式应用于 \(x = (p_{\mathrm{M}} - p_{\mathrm{N}}).y\)，并注意关系式 \((1_{\mathrm{E}} - Q_{\mathrm{M}})(p_{\mathrm{M}} - p_{\mathrm{N}}) = (Q_{\mathrm{M}} - Q_{\mathrm{N}}) p_{\mathrm{N}}\) 与 \((p_{\mathrm{M}} - p_{\mathrm{N}}) Q_{\mathrm{M}} = (1_{\mathrm{E}} - p_{\mathrm{N}})(Q_{\mathrm{M}} - Q_{\mathrm{N}}).\)）
\end{enumerate}

¶ 7) a) 沿用习题 6 的记号，证明：为使 \(\mathbf{M} + \mathbf{N}^{\circ}\) 是闭的，必要且充分的是 \(\mathbf{M} + \mathbf{N}^{\circ} = (\mathbf{M}^{\circ} \cap \mathbf{N})^{\circ}\)。

\begin{enumerate}
\def\labelenumi{\alph{enumi})}
\setcounter{enumi}{1}
\tightlist
\item
  证明下列性质等价：
\end{enumerate}

\(\alpha)\)

\[
\varepsilon (\mathbf {M}, \mathbf {N}) <   1;
\]

β) \(M + N^{\circ}\) 在 E 中是闭的；

\(\gamma)\) 若 \(\tilde{\mathbf{M}}\) 是 \(\mathbf{M} \cap \mathbf{N}^{\circ}\) 在 \(\mathbf{M}\) 中的正交补，而 \((\mathbf{N}^{\circ})^{\sim}\) 是 \(\mathbf{M} \cap \mathbf{N}^{\circ}\) 在 \(\mathbf{N}^{\circ}\) 中的正交补，则 \(\mathbf{E}\) 是 \(\mathbf{M}^{\circ} \cap \mathbf{N}\)、\(\mathbf{M} \cap \mathbf{N}^{\circ}\)、\(\tilde{\mathbf{M}}\) 与 \((\mathbf{N}^{\circ})^{\sim}\) 的直和。此外，若 \(R\) 与 \(S\) 分别是这一分解所对应的、从 \(\mathbf{E}\) 到 \(\mathbf{M}\) 与 \((\mathbf{N}^{\circ})^{\sim}\) 上的投影，则有 \(\| R \| = \| S \| = (1 - \varepsilon^2 (\mathbf{M}, \mathbf{N}))^{1/2}\)。（为看出 \(\alpha\) ) 蕴涵 \(\beta\) )，首先注意：若 \(x \in \tilde{\mathbf{M}}\) 且 \(y \in (\mathrm{N}^\circ)^{\sim}\)，则 \(|\langle x | y \rangle| \leqslant \varepsilon (\mathbf{M}, \mathbf{N}) \| x \| . \| y \|\)；然后，设 \(u = x + y + t\) 是一个元素 \(u \in \mathbf{M} + \mathrm{N}^\circ\) 的分解，其中 \(x \in \mathbf{M}\)、\(y \in (\mathrm{N}^\circ)^{\sim}\) 且 \(t \in \mathbf{M} \cap \mathrm{N}^\circ\)，由此推出 \(\| x \| \leqslant (1 - \varepsilon^2 (\mathbf{M}, \mathbf{N}))^{-1/2} \| u \|, \| y \| \leqslant (1 - \varepsilon^2 (\mathbf{M}, \mathbf{N}))^{-1/2} \| u \|\) 以及 \(\| t \| \leqslant \| u \|\)。为证明 \(\gamma\) ) 蕴涵 \(\alpha\) )，考察分解 \(v = v_1 + v_2\)（对一个 \(v \in \tilde{\mathbf{M}}\)），其中 \(v_1\) 是 \(v\) 到 \(\tilde{\mathbf{N}}\)（即 \(\mathbf{N} \cap \mathbf{M}^\circ\) 在 \(\mathbf{N}\) 中的正交补）上的正交投影；我们有 \(R.v_1 = v\)，且假若有 \(\varepsilon (\mathbf{M}, \mathbf{N}) = 1\)，就会存在一个序列 \((v_n) \in \mathbf{M}\)，使得 \(\| v_n \| = 1\) 且使得 \(\| (v_n)_1\|\) 趋于 0。其次证明限制 \(R_1\)（即 \(R\) 到 \(\tilde{\mathbf{N}}\) 上的限制）是从 \(\tilde{\mathbf{N}}\) 到 \(\tilde{\mathbf{M}}\) 上的一个双射；为计算 \(\| R \|\)，证明 \(\| R_1^{-1} \| \leqslant (1 - \varepsilon^2 (\mathbf{M}, \mathbf{N}))^{1/2}\)。

\begin{enumerate}
\def\labelenumi{\alph{enumi})}
\setcounter{enumi}{2}
\tightlist
\item
  由 b) 推出：若 \(M + N^{\circ}\) 是闭的，则 \(M^{\circ} + N\) 也是闭的。
\end{enumerate}

\begin{enumerate}
\def\labelenumi{\arabic{enumi})}
\setcounter{enumi}{7}
\tightlist
\item
  \begin{enumerate}
  \def\labelenumii{\alph{enumii})}
  \tightlist
  \item
    设 E 是一个希尔伯特空间，T 是一个从 E 到其自身的连续线性映射，满足 \(\|T\|\leqslant1\)。证明关系式 \(T.x=x,\langle T.x|x\rangle=\|x\|\) 与 \(T^{*}.x=x\) 等价，并且 \(1_{E}-T\) 的核与 \(1_{E}-T\) 的像的闭包互为正交补。
  \end{enumerate}
\end{enumerate}

\begin{enumerate}
\def\labelenumi{\alph{enumi})}
\setcounter{enumi}{1}
\tightlist
\item
  设 \(T\) 是一个从 \(E\) 到其自身的连续线性映射，满足不等式

\[
\| x - T. x \| ^ {2} \leqslant \| x \| ^ {2} - \| T. x \| ^ {2}\tag{1}
\]

对所有 \(x \in E\) 成立。于是 \(\|T.x\| < \|x\|\) 对所有使得 \(T.x \neq x\) 的 x 成立；对每个 \(x \in E\)，序列 \((T^{n}.x)\) 收敛于一点 P.x，且 P 是到 \(1_{E} - T\) 的核上的正交投影算子。
\item
  设 \(P_{1}, \ldots, P_{r}\) 是 E 上的正交投影算子。证明乘积 \(T = P_{1}P_{2} \ldots P_{r}\) 满足关系式 (1)（对 r 用归纳法论证）；于是正交投影算子 P 是到诸投影算子 \(P_{j}\) 的像的交上的正交投影算子（注意，若对一个指标 j 有 \(\|P_{j}.x\| < \|x\|\)，则 \(\|T.x\| < \|x\|\)）。
\end{enumerate}

\begin{enumerate}
\def\labelenumi{\arabic{enumi})}
\setcounter{enumi}{8}
\tightlist
\item
  设 E 是一个希尔伯特空间，\((P_{j})_{j\in\mathbb{N}}\) 是 E 上正交投影算子的一个序列，使得对所有 \(j\in N\)，存在一个 \(n_{j}\in N\)，使得对所有 \(k\in N\)，正交投影算子 \(P_{k}, P_{k+1}, ..., P_{k+n_{j}}\) 中至少有一个等于 \(P_{j}\)。置 \(R_{s}=P_{s}P_{s-1}\ldots P_{0}\)（对所有 \(s\in N\)）。
\end{enumerate}

\begin{enumerate}
\def\labelenumi{\alph{enumi})}
\tightlist
\item
  对每个 \(x\in E\)，令 \(x_{s}=R_{s}\cdot x\)。证明 \(\sum_{s}\|x_{s-1}-x_{s}\|^{2}\leqslant\|x\|^{2}\)，并由此推出，对
\end{enumerate}

每个整数 \(r \geqslant 1\)，\(x_{s + r} - x_s\) 当 \(s\) 趋于 \(+\infty\) 时趋于 0。

\begin{enumerate}
\def\labelenumi{\alph{enumi})}
\setcounter{enumi}{1}
\item
  设 \((x_{s_k})\) 是从 \((x_s)\) 中抽取的一个序列，它弱收敛于一个极限 \(y\)；于是每个序列 \((x_{s_k + r})\) 也弱收敛于 \(y\)。由此推出，\(y\) 属于每个子空间 \(\mathbf{M}_j = P_j(\mathbf{E})\)。（对每个 \(j\)，存在 \(r_k\) 使得 \(0 \leqslant r_k \leqslant n_j\) 且 \(s_k + r_k = j\)；证明序列 \((x_{s_k + r_k})\) 弱收敛于 \(y\)。）
\item
  证明序列 \((x_{s})\) 弱收敛于 \(x\) 到交 \(\mathbf{M}\)（诸 \(\mathbf{M}_{j}\) 的交）上的正交投影。（化归到 \(\mathbf{M} = \{0\}\) 的情形，并利用 \(b\) ) 以及 E 中每个闭球的弱紧性。）
\end{enumerate}

¶ 10) 设 E 是一个希尔伯特空间，u 是 E 的一个正算子。

\begin{enumerate}
\def\labelenumi{\alph{enumi})}
\tightlist
\item
  证明：对所有 \(x \in E\)，我们有

\[
\left\| u (x) \right\| ^ {2} \leqslant \| u \|. \langle u (x) | x \rangle
\]

(注意 \(\langle u(x)|u(x)\rangle^2\leqslant \langle u(x)|x\rangle \langle u^2 (x)|u(x)\rangle\)，见 V, p.~3, prop. 2)。
\item
  设 M 是 E 的一个闭向量子空间，\(\mathbf{M}^{\circ}\) 是它的正交补。设 \(x\in \mathbf{M}\)，并令 \(f(x)\) 为 \(\langle u(x + y)|x + y\rangle\)（当 \(y\) 取遍 \(\mathbf{M}^{\circ}\) 时）的下确界。对每个 \(\varepsilon >0\)，令 \(\operatorname {E}(x,\varepsilon)\) 为由所有 \(y\in \mathbf{M}^{\circ}\)（满足 \(\langle u(x + y)|x + y\rangle \leqslant f(x) + \varepsilon\)）组成的集。证明 \(\operatorname {E}(x,\varepsilon)\) 是凸的，并且对所有 \(z\in \mathbf{M}^{\circ}\)，我们有

\[
\langle u (x + y) | z \rangle^ {2} \leqslant \varepsilon \langle u (z) | z \rangle \quad \text { for } \quad y \in \mathrm{E} (x, \varepsilon)\tag{*}
\]

(考察实变量 t 的函数 \(g:t \mapsto \langle u(x + y + tz)/x + y + tz \rangle\)，它在一点 \(t_{0}\) 处达到其极小值，并注意 \(g(t_{0}) \geqslant f(x)\) 与 \(g(0) \leqslant f(x) + \varepsilon\) )。
\item
  对每个整数 \(n \geqslant 1\)，取 \(y_{n} \in \mathrm{E}(x, 1/n)\)；证明序列 \((u(x + y_{n}))\) 趋于一个属于 M 的极限 \(x_{1}\)，且序列 \((\langle u(y_{n})|y_{n}\rangle)\) 有界（利用不等式 (*)，为数 \(\langle u(y_{n} - y_{m})|y_{n} - y_{m}\rangle\)（对 \(m \geqslant n\)）以及数 \(|\langle u(x + y_{n})|z\rangle|\)（对 \(z \in M^{\circ}\)）各找出一个界）。

\item
  设 \((y_n')\) 是 \(\mathbf{M}^\circ\) 中点的一个序列，使得 \(\langle u(y_n' - y_m')|y_n' - y_m'\rangle\) 当 \(m\) 与 \(n\) 充分大时任意小，且使得序列 \((u(x + y_n'))\) 有一个极限 \(x_1' \in \mathbf{M}\)；

证明 \(x_1' = x_1\)。（令 \(Q(z) = \langle u(z)|z \rangle\)；首先证明数 \(Q((y_p - y_p') - (y_q - y_q'))\) 一旦 \(p\) 与 \(q\) 充分大就任意小，并由此推出序列 \((Q(y_n - y_n'))\) 有界；利用 \(\langle x_1' - x_1|y_p - y_p' \rangle = 0\)（对所有 \(p\)）这一事实，证明序列 \((Q(y_n - y_n'))\) 趋于 0，并利用 \(a\) )。

\item
  由 \(d\) ) 推出：点 \(x_{1}\) 不依赖于 \(y_{n} \in \mathrm{E}(x, 1/n)\) 的选取，且若置 \(u_{1}(x) = x_{1}\)，则 \(u_{1}\) 是一个从 \(\mathbf{M}\) 到其自身的线性映射。证明 \(0 \leqslant \langle u_{1}(x)|x\rangle \leqslant \langle u(x)|x\rangle\) 对所有 \(x \in \mathbf{M}\) 成立，因而 \(u_{1}\) 连续，并且是 \(\mathbf{M}\) 的一个 \(\geqslant 0\) 的自同态。（注意 \(\langle u(x + y_n)|y_n\rangle\) 趋于 0 且 \(\langle u(x + y_n)|x + y_n\rangle\) 趋于 \(f(x)\)。）
\item
  设 \(p_{M}\) 是以 M 为像的正交投影算子，并令 \(u_{0} = u_{1} \circ p_{M}\)。我们有 \(0 \leqslant u_{0} \leqslant u\)，M 在 \(u_{0}\) 下稳定，且 \(u_{0}\) 到 \(M^{\circ}\) 上的限制为零。证明 \(u_{0}\) 是由所有满足 \(v \geqslant 0\)、\(v \leqslant u\)、M 在 v 下稳定且 v 到 \(M^{\circ}\) 上的限制为零的自同态 v 组成的族中的最大元素。
\end{enumerate}

\begin{enumerate}
\def\labelenumi{\arabic{enumi})}
\setcounter{enumi}{10}
\tightlist
\item
  设 \(\mathrm{E}\) 与 \(\mathrm{F}\) 是两个希尔伯特空间。证明：对每个元素 \(u\)（\(\mathrm{E} \hat{\otimes}_2 \mathrm{F}\) 中的），存在一个标准正交序列 \((e_n)\)（\(\mathrm{E}\) 中的）、一个标准正交序列 \((f_n)\)（\(\mathrm{F}\) 中的）以及一个序列 \((\lambda_n)\)（其诸数 \(\geqslant 0\)），使得 \(\sum_{n} \lambda_n^2 < +\infty\) 且 \(u = \sum_{n} \lambda_n e_n \otimes f_n\)；于是 \(\| u \|_2^2 = \sum_{n} \lambda_n^2\)
\end{enumerate}

(cf.~V, p.~55, th. 2 及 p.~53, th. 1)。

¶ 12) 设 E 是一个实希尔伯特空间，V 是 E 中一个以 0 为顶点的闭凸锥，令 V° 为 V 的极锥（在 E 中，把 E 典范地与其对偶等同）。

\begin{enumerate}
\def\labelenumi{\alph{enumi})}
\item
  证明每个点 \(x \in \mathbf{E}\) 都可以唯一地写成形式 \(x = x_{+} - x_{-}\)，其中 \(x_{+} \in \mathbf{V}\) 且 \(x_{-} \in \mathbf{V}^{\circ}\)，以及 \(\langle x_{+}|x_{-}\rangle = 0\)。
\item
  对 V 的每个刻面 F（II, p.~87, exerc. 3），F 或是点 0，或是一个以 0 为顶点的凸锥；所有与 F 正交的 \(y \in V^{\circ}\) 组成的集是一个闭刻面 \(F'\)（\(V^{\circ}\) 的）（但它不是 II, p.~87, exerc. 6 意义下 F 的「对偶刻面」，后者是空集）。
\item
  取 E 为一个实希尔伯特空间的全体 Hilbert-Schmidt 自同态组成的集，取 V 为 E 中正元素组成的集。证明我们有 \(V^{\circ} = V\)，并在此情形解释 a) 的结果（为看出 \(V \subset V^{\circ}\)，利用 V, p.~56 的 cor. 1）。
\item
  在 c) 的假设下，设 \(v \in V\)；所有 \(x \in H\)（满足 \(\langle v(x)|x \rangle = 0\)，或等价地 \(v(x) = 0\)，V, p.~77, exerc. 10）组成的集 L 是 H 的一个闭向量子空间，且 v 在 V 中的刻面 F 是由所有这样的 \(u \in V\) 组成的闭集：\(u(x) = 0\) 对所有 \(x \in L\) 成立；它可以与希尔伯特空间 \(L^{\circ}\) 的全体正 Hilbert-Schmidt 自同态组成的锥等同。由此推出：从 E 到凸集 F 上的投影（V, p.~11）与从 E 到由 F 生成的 E 的闭向量子空间上的正交投影算子恒同。
\end{enumerate}

¶ 13) 设 E 是一个希尔伯特空间，G 是 E 的希尔伯特空间结构的全体自同构组成的群的一个子群。令 \(E^{G}\) 为 E 中由在 G 下不变的所有向量组成的闭向量子空间，p 为从 E 到 \(E^{G}\) 上的正交投影算子。

\begin{enumerate}
\def\labelenumi{\alph{enumi})}
\item
  证明 \(\mathbf{E}^{\mathrm{G}}\) 在 \(\mathbf{E}\) 中的正交补是由向量 \(s.x - x\)（其中 \(s\) 取遍 \(\mathbf{G}\) 且 \(x \in \mathbf{E}\)）生成的闭向量子空间。
\item
  设 \(\mathbf{H}\) 是 \(\mathbf{E}\) 的一个在 \(\mathbf{G}\) 下稳定的非空闭凸子集。证明 0 到 \(\mathbf{H}\) 上的投影属于 \(\mathbf{E}^{\mathrm{G}}\)。
\item
  假设 H 是 E 中一点 x 的轨道的闭凸包，令 a 为 0 到 H 上的投影。证明 \(a = p(x)\)，且 \(H \cap E^{G}\) 缩为点 a（「Birkhoff-Alaoglu 定理」）。（注意 x - a 包含在 \(E^{G}\) 的正交补中。）
\item
  假设 G 由 E 的一个自同构 \(u\) 生成。证明 \(p(x) = \lim_{n\to \infty}\frac{1}{n + 1}\sum_{j = 0}^{n}u^{j}(x)\)
\end{enumerate}

对所有 \(x \in \mathrm{E}\) 成立（若 \(y_{n} = \frac{1}{n + 1} \sum_{j=0}^{n} u^{j}(x)\)，注意序列 \((y_{n})\) 有一个弱极限点 \(a\)，且 \(u(a) = a\)，然后利用 \(c\) )。

\begin{enumerate}
\def\labelenumi{\alph{enumi})}
\setcounter{enumi}{4}
\tightlist
\item
  假设 \(G\) 是一个同态 \(t \mapsto u_t\)（从 \(\mathbf{R}\) 到 \(E\) 的自同构群上的）的像，使得对所有 \(x \in E\)，\(t \mapsto u_t \cdot x\) 是一个从 \(\mathbf{R}\) 到 \(E\) 中的连续映射。证明

\(p(x) = \lim_{T \to \infty} \frac{1}{T} \int_{0}^{T} u_t \cdot x \, dt\)

对所有 \(x \in E\) 成立。

\end{enumerate}

\begin{enumerate}
\def\labelenumi{\alph{enumi})}
\setcounter{enumi}{5}
\tightlist
\item
  假设存在一个元素 \(x \neq 0\)（属于 \(\mathbf{E}\)）和一个数 \(\alpha\)，使得 \(0 < \alpha < 1\) 且 \(\| s.x - x\| \leqslant \alpha \| x\|\) 对所有 \(s \in \mathbf{G}\) 成立。证明 \(\mathrm{E}^{\mathrm{G}} \neq \{0\}\)（利用 \(c\)）。
\end{enumerate}

\begin{enumerate}
\def\labelenumi{\arabic{enumi})}
\setcounter{enumi}{13}
\tightlist
\item
  设 E 是一个复希尔伯特空间，T 是 E 的一个 Hilbert-Schmidt 自同态。a) 设 R 与 L 是正 Hilbert-Schmidt 自同态，满足 \(R^{2} = T^{*}T\) 与 \(L^{2} = TT^{*}\)（V, p.~57, cor. 3）；令 \(R = \text{abs}(T)\)，并称之为 T 的「绝对值」（cf.~V, p.~75, exerc. 4）；我们有 \(L = \text{abs}(T^{*})\)。证明 \(\text{Ker}(T) = \text{Ker}(R)\) 且 \(\overline{L(\text{E})} = \overline{T(\text{E})}\)。存在唯一的等距 V，它从 \(R(\text{E})\) 到 \(T(\text{E})\) 上，使得 T = VR；若把 V 连续地延拓到 \(\overline{R(\text{E})}\)，再延拓为一个算子 \(U \in \mathcal{L}(\text{E})\)（通过在 \(R(\text{E})\) 的正交补上取 U.x = 0），则我们仍有 T = UR（T 的极分解）。于是 \(R = U^{*}T = U^{*}UR = RU^{*}U\) 且 \(L = URU^{*}\)、\(T = LU^{*}\)。若 T 属于 \(\mathcal{L}^{1}(\text{E})\)，则 \(R = \text{abs}(T)\) 亦然，且 T 是两个 Hilbert-Schmidt 自同态的乘积。
\end{enumerate}

\begin{enumerate}
\def\labelenumi{\alph{enumi})}
\setcounter{enumi}{1}
\item
  若 \(T\) 属于 \(\mathcal{L}^1 (\mathrm{E})\)，证明 \(\operatorname {Tr}(\operatorname {abs}(T)) = \sup (\sum_i |\langle a_i | T.b_i\rangle |)\)，其中在右端，\((a_{i})\) 与 \((b_{i})\) 取遍 E 的标准正交基组成的集（利用 \(T\) 的极分解）。证明：若置 \(\| T\| _1 = \operatorname {Tr}(\operatorname {abs}(T))\)，则 \(\| T\| _1\) 是空间 \(\mathcal{L}^1 (\mathrm{E})\) 上的一个范数，使得 \(\| T\| _2\leqslant \| T\| _1\)。
\item
  反之，若 \(T \in \mathcal{L}(\mathrm{E})\) 使得对每一对 \(((a_i), (b_i))\)（\(\mathbf{E}\) 的标准正交基的），和式 \(\sum_{i} |\langle a_i | T \cdot b_i \rangle|\) 都有限，则 \(T \in \mathcal{L}^1(\mathrm{E})\)（首先注意 \(T\) 是一个 Hilbert-Schmidt 自同态，然后利用 \(T\) 的极分解）。
\item
  设 \((\bar{T}_{\nu})\) 是 Hilbert-Schmidt 自同态（相应地 \(\mathcal{L}^1 (\mathrm{E})\) 的元素）的一个序列，使得对 E 的每对点 \(x,y\)，序列 \((\langle x|T_{\nu}.y\rangle)\) 收敛于 \(\langle x|T.y\rangle\)，其中 \(T\) 是一个从 E 到其自身的线性映射；此外，假设范数序列 \(\| T_{\nu}\| _2\)（相应地 \(\| T_{\nu}\| _1\)）有界。证明 \(T\) 是一个 Hilbert-Schmidt 自同态（相应地 \(\mathcal{L}^1 (\mathrm{E})\) 的一个元素）（利用 \(b\) ) 与 \(c\) )）。
\item
  由 \(d\) ) 推出：空间 \(\mathcal{L}^1 (\mathrm{E})\) 对范数 \(\| T\| _1\) 是一个 Banach 空间。
\item
  为使一个自同态 \(T \in \mathcal{L}(\mathrm{E})\) 属于 \(\mathcal{L}^1(\mathrm{E})\)，必要且充分的是：对至少一个标准正交基 \((e_i)\)（\(\mathrm{E}\) 的），和式 \(\sum_{i} \| T.e_i\|\) 有限（沿用 \(a\) ) 的记号，注意 \(|\langle e_i|R.e_i\rangle| \leqslant \| T.e_i\|\)）。

\item
  在空间 \(\ell_{\mathbf{C}}^2\) 中，令 \((e_n)\) 为典范标准正交基，并令 \(a = \sum_{n=0}^{\infty} \frac{1}{n+1} e_n\)；

若 \(\mathrm{F}\) 是一维子空间 \(\mathbf{C}.a\)，则正交投影算子 \(p_{\mathrm{F}}\) 有有限迹，但级数 \(\sum_{n}\| p_{\mathrm{F}}\cdot e_n\|\) 不收敛。
\end{enumerate}

\begin{enumerate}
\def\labelenumi{\arabic{enumi})}
\setcounter{enumi}{14}
\tightlist
\item
  设 \(\mathbf{E}\) 是一个复希尔伯特空间；令 \(\mathcal{B} = \mathcal{L}(\mathbf{E})\) 表示 \(\mathbf{E}\) 的连续自同态组成的代数，赋以通常的范数 \(\| T\| = \sup_{\| x\| \leqslant 1}\| T.x\|\)。对 E 的每对点 \(x, y\)，令 \(\omega_{x,y}\) 表示连续线性型 \(T \mapsto \langle x|T.y \rangle\)（\(\mathcal{B}\) 上的），并令 \(\mathcal{B}_{\circ}\) 为强对偶 \(\mathcal{B}'\)（Banach 空间 \(\mathcal{B}\) 的）中由诸 \(\omega_{x,y}\) 生成的闭子空间。a) 证明：把每个 \(T \in \mathcal{B}\) 对应到线性型 \(\omega \mapsto \langle \omega, T \rangle\)（\(\mathcal{B}_{\circ}\) 上的）的线性映射，是从 \(\mathcal{B}\) 到 \(\mathcal{B}_{\circ}\) 的强对偶上的一个等距；换句话说，\(\mathcal{B}_{\circ}\) 是 \(\mathcal{B}\) 的一个预对偶（IV, p.~56, exerc. 23）。
\end{enumerate}

拓扑 \(\sigma (\mathcal{B},\mathcal{B}_{\circ})\)（\(\mathcal{B}\) 上的）称为超弱拓扑。

\begin{enumerate}
\def\labelenumi{\alph{enumi})}
\setcounter{enumi}{1}
\item
  对每个元素 \(T\)（\(\mathcal{L}^1(\mathrm{E})\) 中的），我们定义一个线性型 \(\phi_T\)（\(\mathcal{B}\) 上的），所用公式为 \(\phi_T(S) = \operatorname{Tr}(ST)\)（对每个算子 \(S \in \mathcal{B}\)）。证明 \(\phi_T\) 连续，且映射 \(T \mapsto \phi_T\) 是从 Banach 空间 \(\mathcal{L}^1(\mathrm{E})\)（exerc. 14, e)）到 Banach 空间 \(\mathcal{B}_{\circ}\) 上的一个等距（先考虑 \(T\) 为有限秩的情形）。
\item
  令 \(\mathcal{B}_{\circ \circ}\) 为 \(\mathcal{B}_{\circ}\) 中由诸 \(\omega_{x,y}\) 生成的向量子空间（于是 \(\mathcal{B}_{\circ} = \overline{\mathcal{B}}_{\circ \circ} \)）。证明 \(\mathcal{B}_{\circ \circ}\) 是桶型的（注意 \(\mathcal{B}\) 的一个子集若对 \(\sigma (\mathcal{B},\mathcal{B}_{\circ \circ})\) 有界，则对范数拓扑有界）。
\item
  令 \(F_{n}\) 为 \(\mathcal{B}_{\circ \circ}\) 的一个子空间，它是 E 的秩 \(\leqslant n\) 的自同态组成的集在 b) 中定义的等距下的像。证明 \(F_{n}\) 在 \(\mathcal{B}_{\circ \circ}\) 中无处稠密，并由此推出 \(\mathcal{B}_{\circ \circ}\) 不是 Baire 空间。
\end{enumerate}

\backmatter
